\section{Appendix: Weighted approximation and constrained moment priors}\label{sec:appendix-weighted}

The observable upper bound in \cref{obj:thm:observable-upper} combines polynomial approximation with factorial estimation. The complementary lower-bound argument requires probability priors whose moments agree while their target values remain separated. This appendix presents the \leanref{S-4}{scalar approximation problem} underlying that construction, the localization bounds used to obtain the priors, and the resulting comparison between approximation error and constrained prior separation.

\subsection{The scalar functional and probability constraints}

The scalar target combines a change in slope at intensity one with the overlap-dependent weighting. We introduce its formula before defining the approximation problem.



The weighted scalar positive-part functional \leanref{sym:\phi_\epsilon(z)}{\ensuremath{\phi_\epsilon(z)}} is
\[
\phi_\epsilon(z)
=\frac{\{(1-\epsilon)+\epsilon z\}\max\{z-1,0\}}{1+z}.
\]
It is continuous and vanishes below intensity one. Its change in slope supplies the local approximation difficulty. The polynomial degree \leanref{sym:K}{\ensuremath{K}} and support endpoint \leanref{sym:M}{\ensuremath{M}} determine the approximation scale. Weighting the uniform error by \(1+z\), where \leanref{sym:z}{\ensuremath{z}} is the scalar intensity, also aligns the approximation criterion with the common-mean restriction on the priors.

\begin{definitionv}{def:weighted-error}[Weighted approximation error \(\mathcal E_{K,M,\epsilon}^{w}\)]
The weighted best uniform polynomial approximation error is
\[
\mathcal E_{K,M,\epsilon}^{w}
:=
\inf_{P\in\mathbb R_K[z]}
\sup_{0\le z\le M}
\frac{|\phi_\epsilon(z)-P(z)|}{1+z},
\]
where \(\phi_\epsilon\) is the weighted scalar positive-part functional and \(\mathbb R_K[z]\) is the set of real polynomials in \(z\) of degree at most \(K\).
\end{definitionv}

The weighted approximation error \leanref{sym:\mathcal E_{K,M,\epsilon}^{w}}{\ensuremath{\mathcal E_{K,M,\epsilon}^{w}}} measures the discrepancy after allowing the pointwise error to grow with intensity. For probability measures of mean one, the weight \(1+z\) has integral two. This normalization gives approximation error a direct role in bounding target differences between moment-matched priors. The admissible probability pairs are defined next.

\begin{definitionv}{def:constrained-prior-separation}[Constrained prior separation \(\Delta_{K,M,\epsilon}\)]
The normalized constrained prior separation is
\[
\Delta_{K,M,\epsilon}
=
\sup_{(\nu_0,\nu_1)\in\mathfrak N_{K,M}}
\left|
\int_0^M\phi_\epsilon(z)\,d(\nu_1-\nu_0)(z)
\right|,
\]
where \(\phi_\epsilon\) is the weighted scalar positive-part functional and the parameters satisfy
\begin{itemize}
\item \textbf{(Degree.)} \(K\) is an integer with \(K\geq2\).
\item \textbf{(Support endpoint.)} \(M\) is real with \(2\leq M\leq K^2\).
\item \textbf{(Overlap.)} \(0\leq\epsilon\leq1/2\).
\end{itemize}
The set \(\mathfrak N_{K,M}\) consists of pairs \((\nu_0,\nu_1)\) of probability measures on \([0,M]\) satisfying
\[
\int_0^M z\,d\nu_0(z)
=
\int_0^M z\,d\nu_1(z)
=1
\]
and
\[
\int_0^M z^j\,d\nu_0(z)
=
\int_0^M z^j\,d\nu_1(z),
\qquad j=0,1,\ldots,K.
\]
\end{definitionv}

The constrained prior separation \leanref{sym:\Delta_{K,M,\epsilon}}{\ensuremath{\Delta_{K,M,\epsilon}}} therefore incorporates both moment agreement and an absolute normalization of the first moment. Moment agreement cancels polynomial target components through the prescribed degree, while the common mean fixes the average intensity. The parameter range in \cref{obj:def:constrained-prior-separation} is the range covered by the comparison below, including zero overlap for this scalar approximation problem.

A signed measure supplies the target contrast; a common atom completes its two nonnegative parts into probability priors. We define that completion as a map on measures. The support and normalization properties needed for the constructed packet inputs are established in \cref{obj:thm:weighted-separation-and-frontier}.

\begin{definitionv}{def:weighted-handle}[Common-atom completion map \(\mathscr H^w\)]
The common-atom completion map is
\[
\mathscr H^w(\sigma_+,\sigma_-)
=
\left(
\sigma_-+(1-t)_+\delta_{z_0},
\ \sigma_++(1-t)_+\delta_{z_0}
\right),
\]
with
\[
t=\sigma_+(\mathbb R),
\qquad
u=\int_{\mathbb R}z\,d\sigma_+(z),
\qquad
z_0=\frac{1-u}{1-t}.
\]
Its inputs satisfy the following conditions:
\begin{itemize}
\item \textbf{(Measures.)} \(\sigma_+\) and \(\sigma_-\) are finite nonnegative measures on \(\mathbb R\).
\item \textbf{(First moments.)} Each measure has a finite first moment.
\item \textbf{(Mass.)} \(t<1\).
\end{itemize}
Here \((r)_+=\max\{r,0\}\), and \(\delta_{z_0}\) is the unit point mass at \(z_0\).
\end{definitionv}

The common-atom completion map \leanref{sym:\mathscr H^{w}}{\ensuremath{\mathscr H^w}} adds the same measure to both inputs, preserving their signed difference. Its location uses the positive input's remaining mass and first moment. For the packet inputs in \cref{obj:thm:weighted-separation-and-frontier}, the stated bounds place this atom inside the intensity support and certify that both completed measures have total mass and mean one.

\subsection{Fourier coefficients and packet localization}

A modulated Jackson function provides a localized signed contrast with a prescribed frequency band. Its twice antiderivative separates the response at the change in slope from the contribution spread across the interval. The following grouped notation fixes the periodic normalization, coefficient conventions, and endpoint extension before stating the localization bounds.



The shifted squared frequencies in \cref{obj:lem:jackson-packet-localization} encode two integrations of the modulated function. The endpoint calculation concerns the coefficients after extension by zero: forward differences crossing either end of the coefficient band enter the same sum as interior differences. This convention gives a bound for the complete Fourier expression. The localization statement records that coefficient bound together with frequency cancellation, the value at the periodic origin, and spatial decay.

\begin{lemmav}{lem:jackson-packet-localization}[Jackson packet localization bounds]
There exist real constants \(0<c\leq C\) such that the following holds for every integer \(N\geq2\). Define the normalized periodic Jackson function by
\[
\mathcal J_N(u)
=c_N\left\{\frac{\sin(Nu/2)}{\sin(u/2)}\right\}^{4},
\qquad
c_N=
\left(
\int_{-\pi}^{\pi}
\left\{\frac{\sin(Nu/2)}{\sin(u/2)}\right\}^{4}
\frac{du}{2\pi}
\right)^{-1},
\]
with removable values filled continuously. Set
\[
L=4N,\qquad
k_N(u)=\mathcal J_N(u)\cos(Lu),
\qquad
\widehat h(j)=\int_{-\pi}^{\pi}h(u)e^{-iju}\,\frac{du}{2\pi}.
\]
For \(j\in\mathbb Z\), define the zero-extended coefficient sequences and their forward differences by
\[
a_j^\pm=
\begin{cases}
\widehat{\mathcal J_N}(j)/(L\pm j)^2,& |j|\leq2N-2,\\
0,& |j|>2N-2,
\end{cases}
\qquad
(\delta a)_j=a_{j+1}-a_j,
\qquad
(\delta^2a)_j=a_{j+2}-2a_{j+1}+a_j.
\]
Define also
\[
G_N(u)
=-\frac12\sum_{j=-(2N-2)}^{2N-2}
\left\{
a_j^+\cos((L+j)u)+a_j^-\cos((L-j)u)
\right\},
\qquad
\|u\|_{\mathbb T}
=\inf_{j\in\mathbb Z}|u-2\pi j|.
\]
Then
\[
\int_{-\pi}^{\pi}\mathcal J_N(u)\,\frac{du}{2\pi}=1.
\]
For each sign, the second forward differences are absolutely summable and satisfy
\[
\sum_{j\in\mathbb Z}|\delta^2a_j^\pm|
\leq\frac{C}{N^3}.
\]
The differences include the endpoints of the displayed zero extension. For every \(r\in\mathbb Z\) satisfying
\[
|r|<2N+2\quad\text{or}\quad |r|>6N-2,
\]
the trigonometric integrals vanish:
\[
\int_{-\pi}^{\pi}k_N(u)\cos(ru)\,du=0,
\qquad
\int_{-\pi}^{\pi}k_N(u)\sin(ru)\,du=0.
\]
The function \(G_N\) satisfies
\[
G_N''(u)=k_N(u)\quad\text{for every }u\in\mathbb R,
\qquad
\int_{-\pi}^{\pi}G_N(u)\,du=0,
\]
and the bounds
\[
\frac{c}{N}\leq |G_N(0)|\leq\frac{C}{N},
\qquad
|G_N(u)|\leq
\frac{C}{N\{1+N\|u\|_{\mathbb T}\}^{2}}
\quad\text{for every }u\in\mathbb R,
\qquad
\int_{-\pi}^{\pi}|G_N(u)|\,\frac{du}{2\pi}
\leq\frac{C}{N^2}.
\]
\end{lemmav}

Frequency cancellation and spatial localization play complementary roles in \cref{obj:lem:jackson-packet-localization}. The frequency band supplies the cancellation required for matching low-order moments after the packet is transferred to the intensity interval. The twice antiderivative has an origin value of order \(N^{-1}\) and an integrated absolute magnitude of order \(N^{-2}\). These two scales distinguish the localized response at the scalar target's change in slope from the contribution distributed away from that point.

The second-difference bound includes the Fourier endpoints explicitly. Together with the spatial bound, it controls the entire periodic twice antiderivative, including the tails measured by distance from the periodic origin. The normalized packet used below combines this localization with a calibration that leaves sufficient mass and first moment for the common-atom completion.

\subsection{Weighted separation and statistical consequences}

The packet localization lemma supplies low-frequency cancellation and spatial concentration, and the weighted-separation theorem converts the packet into common-mean moment priors and an approximation lower bound. \Cref{sec:appendix-lower-bound} uses those priors in the statistical comparison. The detailed approximation, sparse, dense, transfer, and concluding arguments are divided explicitly in \cref{sec:deferred-proofs}.

The next result collects the approximation comparison and the completed-prior construction. It also records their statistical consequences alongside the attaining upper bound. Within the scalar parameter range, weighted approximation error and common-mean moment-prior separation have the same order, uniformly over the overlap parameter.

\begin{theoremv}{thm:weighted-separation-and-frontier}[Weighted separation and minimax frontier]
There exist universal real constants \(0<c<C<\infty\) and an integer \(A_0\geq1\) such that the following approximation and separation guarantees hold for every choice of parameters satisfying:
\begin{itemize}
\item \textbf{(Degree.)} \(K\) is an integer with \(K\geq2\).
\item \textbf{(Support.)} \(M\) is real and \(2\leq M\leq K^2\).
\item \textbf{(Overlap.)} \(0\leq\epsilon\leq1/2\).
\end{itemize}
For the weighted error and constrained-prior separation in
\cref{obj:def:weighted-error,obj:def:constrained-prior-separation}, with scalar functional
\[
\phi_\epsilon(z)=
\frac{((1-\epsilon)+\epsilon z)\max\{z-1,0\}}{1+z},
\]
one has
\[
c\frac{\sqrt M}{K}
\leq \mathcal E_{K,M,\epsilon}^{w}
\leq \Delta_{K,M,\epsilon}
\leq 4\mathcal E_{K,M,\epsilon}^{w}
\leq C\frac{\sqrt M}{K}.
\]
Moreover, let \(\sigma_+\) and \(\sigma_-\) be the positive and negative Jordan parts of the normalized Jackson packet with calibration \(A_0\), pushed to \([0,M]\). These are finite measures with integrable first moments. Their positive mass and common completion location satisfy
\[
t=\sigma_+(\mathbb R)\leq\frac12,
\qquad
u=\int z\,d\sigma_+(z),
\qquad
z_0=\frac{1-u}{1-t},
\qquad
\frac12\leq z_0\leq2.
\]
There exists \((\nu_0,\nu_1)\in\mathfrak N_{K,M}\), the constrained prior class of \cref{obj:def:constrained-prior-separation}, such that the common-atom completion of \cref{obj:def:weighted-handle} satisfies
\[
\mathscr H^w(\sigma_+,\sigma_-)
=(\nu_0,\nu_1)
=\bigl(\sigma_-+(1-t)\delta_{z_0},
       \sigma_++(1-t)\delta_{z_0}\bigr),
\]
and
\[
c\frac{\sqrt M}{K}
\leq
\left|
\int\phi_\epsilon(z)\,d\nu_1(z)
-
\int\phi_\epsilon(z)\,d\nu_0(z)
\right|.
\]

There also exist universal real constants
\(0<c_\star\leq C_\star<\infty\), \(H_0>0\), and \(\kappa>0\), and a universal integer \(D_0\geq2\), such that the armwise estimator
\(\widehat V_{n,d,\epsilon}^{\mathrm{AF}}\) with these calibrations satisfies the following risk guarantees for every choice of parameters satisfying:
\begin{itemize}
\item \textbf{(Sample size.)} \(n\) is an integer with \(n\geq1\).
\item \textbf{(Alphabet.)} \(d\) is an integer with \(d\geq2\).
\item \textbf{(Overlap.)} \(0<\epsilon\leq1/2\).
\end{itemize}
The observed risks obey
\[
c_\star\min\left\{1,\frac{d}{n\epsilon\log(ed)}\right\}
\leq \mathfrak R_{n,d,\epsilon}
\leq
\sup_{\mathbb P\in\mathcal D_{d,\epsilon}^{\mathrm{obs}}}
E_{\mathbb P}\!\left[
\left\{\widehat V_{n,d,\epsilon}^{\mathrm{AF}}
-\Psi(\mathbb P)\right\}^2
\right]
\leq
C_\star\min\left\{1,\frac{d}{n\epsilon\log(ed)}\right\}.
\]
The causal minimax risk obeys
\[
c_\star\min\left\{1,\frac{d}{n\epsilon\log(ed)}\right\}
\leq
\inf_{\widehat V}
\sup_{P\in\mathcal P_{d,\epsilon}}
E_{\operatorname{Obs}(P)}\!\left[
\left\{\widehat V-V^\star(P)\right\}^2
\right]
\leq
C_\star\min\left\{1,\frac{d}{n\epsilon\log(ed)}\right\},
\]
where the infimum ranges over measurable real-valued estimators based on \(n\) observed units.

Finally, consider any sequences satisfying:
\begin{itemize}
\item \textbf{(Sample sizes.)} The integer sample sizes \(n_k\) tend to infinity.
\item \textbf{(Alphabets.)} Each \(d_k\) is an integer with \(d_k\geq2\).
\item \textbf{(Overlap.)} \(0<\epsilon_k\leq1/2\) for every \(k\).
\end{itemize}
There exists a sequence of measurable real-valued estimators
\(\widehat V_k\), each based on \(n_k\) observed units with alphabet size \(d_k\), such that
\[
\sup_{\mathbb P\in\mathcal D_{d_k,\epsilon_k}^{\mathrm{obs}}}
E_{\mathbb P}\!\left[
\left\{\widehat V_k-\Psi(\mathbb P)\right\}^2
\right]\longrightarrow0
\]
if and only if
\[
\frac{d_k}{n_k\epsilon_k\log(ed_k)}\longrightarrow0.
\]
\end{theoremv}

The approximation comparison expresses the same difficulty in two forms. Polynomial approximation measures how accurately the scalar target can be represented through low-order moments; constrained separation measures how far target values can differ when those moments agree. The packet construction realizes a separation of order \(\sqrt M/K\) using probability priors with common mean one. Adding the shared completion atom preserves the target contrast, while the bounds on its location place the completed priors within the prescribed support.

The factor four in the upper comparison reflects the normalization in \cref{obj:def:constrained-prior-separation}: each prior assigns integral two to the approximation weight \(1+z\). The remaining inequalities in \cref{obj:thm:weighted-separation-and-frontier} establish both the reverse comparison and the uniform approximation scale. In particular, the explicit completed pair achieves the stated lower order under the probability and mean constraints.

The statistical clauses connect this scalar construction to the main-body results. The constrained priors supply the separation used by the all-estimator lower bound in \cref{obj:thm:all-estimator-lower}, while \cref{obj:thm:observable-upper} supplies the attaining risk bound for the estimator of \cref{obj:def:armwise-estimator}. Their combination gives the matched observed and causal rates stated in \cref{obj:thm:matched-frontier}, with the causal interpretation supported by \cref{obj:prop:identification}. Along the admissible sequences, the same two-sided risk comparison yields the necessary and sufficient condition for uniform consistency in \cref{obj:thm:consistency-threshold}.
